\chapter{Analyse et Architecture}

\section{Introduction}

Ce chapitre formalise les besoins de VM Marketplace, identifie les acteurs du système,
établit le backlog produit et décrit l'architecture technique ainsi que les choix
technologiques qui sous-tendent l'implémentation.

% ─────────────────────────────────────────────────────────────────────────────
\section{Identification des Acteurs}
% ─────────────────────────────────────────────────────────────────────────────

La plateforme distingue deux acteurs :

\begin{description}
  \item[Visiteur (utilisateur non authentifié)] Toute personne accédant à l'application
        web VM Marketplace sans s'être authentifiée. Un visiteur peut accéder au moteur
        de recommandation IA et à la page marketplace pour explorer les modèles de VM,
        mais ne peut effectuer aucune opération de provisionnement ni de gestion.

  \item[Utilisateur authentifié (Chercheur / Étudiant)] Toute personne ayant créé un
        compte et s'étant connectée. Un utilisateur authentifié peut provisionner des
        machines virtuelles et des clusters Kubernetes, gérer ses ressources via le
        tableau de bord, accéder aux terminaux SSH, consulter les tableaux de bord de
        supervision et télécharger les fichiers kubeconfig. Toutes les ressources
        provisionnées sont isolées par compte utilisateur.
\end{description}

% ─────────────────────────────────────────────────────────────────────────────
\section{Besoins Fonctionnels}
% ─────────────────────────────────────────────────────────────────────────────

\subsection{Module Authentification}

\begin{itemize}
  \item En tant que visiteur, je veux m'inscrire avec un e-mail et un mot de passe
        afin d'accéder aux fonctionnalités de provisionnement de la plateforme.
  \item En tant que visiteur, je veux me connecter avec mes identifiants afin de
        recevoir un jeton JWT valide pendant 30 minutes.
  \item En tant qu'utilisateur authentifié, je veux que ma session soit sécurisée par
        un jeton Bearer JWT afin que mes ressources restent privées.
\end{itemize}

\subsection{Module Provisionnement de VM}

\begin{itemize}
  \item En tant qu'utilisateur authentifié, je veux sélectionner une configuration de
        VM via un assistant en plusieurs étapes (type d'instance, OS, stockage, réseau,
        région, révision) afin de personnaliser ma machine avant de valider le paiement.
  \item En tant qu'utilisateur authentifié, je veux payer ma VM via un formulaire
        Stripe afin que le coût soit débité de ma carte avant le début du
        provisionnement.
  \item En tant qu'utilisateur authentifié, je veux voir une progression en temps réel
        étape par étape (via Server-Sent Events) pendant le provisionnement afin de
        connaître l'état actuel sans rafraîchir la page.
  \item En tant qu'utilisateur authentifié, je veux recevoir un e-mail de confirmation
        une fois ma VM provisionnée afin de disposer des informations d'accès dans un
        format durable.
\end{itemize}

\subsection{Module Gestion des VM}

\begin{itemize}
  \item En tant qu'utilisateur authentifié, je veux consulter toutes mes VM
        provisionnées dans un tableau de bord afin de les gérer depuis un seul endroit.
  \item En tant qu'utilisateur authentifié, je veux supprimer une VM depuis le tableau
        de bord afin que les ressources HCS correspondantes soient détruites via
        \texttt{terraform destroy}.
  \item En tant qu'utilisateur authentifié, je veux ouvrir un terminal SSH dans le
        navigateur vers ma VM afin d'effectuer des opérations d'administration sans
        installer de client SSH.
  \item En tant qu'utilisateur authentifié, je veux consulter le tableau de bord de
        supervision Netdata de ma VM (CPU, mémoire, disque, réseau) intégré dans la
        plateforme afin d'évaluer l'utilisation des ressources en un coup d'œil.
\end{itemize}

\subsection{Module Recommandation IA}

\begin{itemize}
  \item En tant que visiteur ou utilisateur authentifié, je veux décrire ma charge de
        travail en langage naturel afin que le moteur IA recommande le flavor et la
        configuration de VM appropriés.
  \item En tant que visiteur ou utilisateur authentifié, je veux voir le niveau de
        confiance et le raisonnement de l'IA avec la recommandation afin d'évaluer
        si elle correspond à mes besoins.
\end{itemize}

\subsection{Module Cluster Kubernetes}

\begin{itemize}
  \item En tant qu'utilisateur authentifié, je veux configurer un cluster Kubernetes
        (flavor maître, flavor worker, nombre de workers, taille disque, région) via un
        assistant afin de déployer des clusters multi-nœuds sans expertise Kubernetes.
  \item En tant qu'utilisateur authentifié, je veux payer le cluster via Stripe puis
        suivre la progression du provisionnement étape par étape afin de savoir quand
        le cluster est prêt.
  \item En tant qu'utilisateur authentifié, je veux télécharger le fichier kubeconfig
        de mon cluster depuis le tableau de bord afin de me connecter avec
        \texttt{kubectl} depuis ma machine locale.
  \item En tant qu'utilisateur authentifié, je veux supprimer un cluster depuis le
        tableau de bord afin que toutes les ressources HCS associées soient détruites
        proprement.
\end{itemize}

% ─────────────────────────────────────────────────────────────────────────────
\section{Besoins Non Fonctionnels}
% ─────────────────────────────────────────────────────────────────────────────

\begin{table}[H]
\centering
\caption{Besoins non fonctionnels}
\label{tab:nfr}
\begin{tabularx}{\linewidth}{llX}
\toprule
\textbf{Catégorie} & \textbf{Besoin} & \textbf{Cible} \\
\midrule
Sécurité & Authentification JWT & Jetons expirant après 30 min ; hachage bcrypt des mots de passe \\
Sécurité & Paiement & Données de carte gérées exclusivement par l'API PCI Stripe \\
Sécurité & Isolation SSH & Chaque cluster utilise une paire de clés RSA unique générée par session \\
Performance & Temps de réponse API & Tous les endpoints hors provisionnement répondent en moins de 500 ms \\
Performance & Recommandation IA & Recommandation retournée en moins de 60 secondes \\
Ergonomie & UX de l'assistant & Formulaire multi-étapes avec validation à chaque étape ; aucune donnée perdue en navigation arrière \\
Fiabilité & Suivi de progression & Le flux SSE maintient la connexion pendant toute la durée du provisionnement (5--10 min) \\
Scalabilité & Utilisateurs concurrents & Chaque provisionnement s'exécute dans un thread en arrière-plan ; l'API reste réactive \\
Maintenabilité & Isolation IaC & Chaque VM/cluster possède son propre répertoire \texttt{terraform\_states/} \\
\bottomrule
\end{tabularx}
\end{table}

% ─────────────────────────────────────────────────────────────────────────────
\section{Diagramme de Cas d'Utilisation Global}
% ─────────────────────────────────────────────────────────────────────────────

La figure~\ref{fig:use_case} présente le diagramme de cas d'utilisation global de
VM Marketplace.

\begin{figure}[H]
\centering
\begin{tikzpicture}[
  uc/.style={ellipse, draw, minimum width=3.5cm, minimum height=0.7cm, align=center,
             font=\scriptsize},
  arr/.style={-{Stealth}, dashed, thin},
]
  % Acteurs
  \node[font=\small\bfseries] (visitor) at (-5, 0) {Visiteur};
  \node[font=\small\bfseries, align=center] (user) at (-5,-5)
    {Utilisateur\\Authentifié};

  % Frontière système
  \draw[thick] (-1,-8) rectangle (9,2);
  \node[font=\small, anchor=north west] at (-1,2) {VM Marketplace};

  % Cas d'utilisation
  \node[uc] (reg)   at (4, 1.2) {S'inscrire / Se connecter};
  \node[uc] (ai)    at (4, 0.0) {Recommandation IA};
  \node[uc] (mkt)   at (4,-1.2) {Parcourir la Marketplace};
  \node[uc] (prov)  at (4,-2.4) {Provisionner une VM};
  \node[uc] (dash)  at (4,-3.6) {Gérer le Tableau de Bord};
  \node[uc] (ssh)   at (4,-4.8) {Terminal SSH Navigateur};
  \node[uc] (mon)   at (4,-6.0) {Consulter la Supervision};
  \node[uc] (k8s)   at (4,-7.2) {Provisionner Cluster K8s};

  % Associations Visiteur
  \draw (visitor) -- (reg);
  \draw (visitor) -- (ai);
  \draw (visitor) -- (mkt);

  % Associations Utilisateur
  \draw (user) -- (ai);
  \draw (user) -- (mkt);
  \draw (user) -- (prov);
  \draw (user) -- (dash);
  \draw (user) -- (ssh);
  \draw (user) -- (mon);
  \draw (user) -- (k8s);

  % Include paiement
  \node[uc, fill=yellow!20] (pay)  at (7.5,-2.4) {Paiement Stripe};
  \node[uc, fill=yellow!20] (payk) at (7.5,-7.2) {Paiement Stripe};
  \draw[arr] (prov) -- node[above,font=\tiny]{$\langle\langle$include$\rangle\rangle$} (pay);
  \draw[arr] (k8s)  -- node[above,font=\tiny]{$\langle\langle$include$\rangle\rangle$} (payk);
\end{tikzpicture}
\caption{Diagramme de cas d'utilisation global --- VM Marketplace}
\label{fig:use_case}
\end{figure}

% ─────────────────────────────────────────────────────────────────────────────
\section{Diagramme de Classes}
% ─────────────────────────────────────────────────────────────────────────────

Le modèle de données central est composé de trois entités ORM SQLAlchemy stockées
dans une base de données SQLite, comme illustré à la figure~\ref{fig:class_diagram}.

\begin{figure}[H]
\centering
\begin{tikzpicture}[
  class/.style={rectangle, draw, minimum width=5cm, font=\small},
  arr/.style={-{Stealth[open]}, thick},
]
  % Classe Utilisateur
  \node[class] (user_title) at (0, 0) {\textbf{Utilisateur}};
  \node[class, anchor=north, text width=5cm] (user_attrs) at (user_title.south)
    {\underline{id : Integer (PK)} \\ email : String \\ mot\_de\_passe\_hash : String \\
     date\_creation : DateTime};

  % Classe MachineVirtuelle
  \node[class] (vm_title) at (7.5, 0) {\textbf{MachineVirtuelle}};
  \node[class, anchor=north, text width=7.5cm] (vm_attrs) at (vm_title.south)
    {\underline{id : Integer (PK)} \\ utilisateur\_id : Integer (FK) \\
     nom\_instance, id\_vm\_cloud : String \\ statut : String \\
     id\_flavor, id\_image : String \\ ip\_publique, ip\_privee : String \\
     url\_netdata : String \\ date\_creation : DateTime};

  % Classe Cluster
  \node[class] (cl_title) at (7.5, -6) {\textbf{Cluster}};
  \node[class, anchor=north, text width=7.5cm] (cl_attrs) at (cl_title.south)
    {\underline{id : Integer (PK)} \\ utilisateur\_id : Integer (FK) \\
     nom, statut : String \\ nb\_workers : Integer \\
     ip\_publique\_maitre, ip\_privee\_maitre : String \\
     ips\_workers : Text (JSON) \\ kubeconfig : Text \\
     date\_creation : DateTime};

  % Associations
  \draw[arr] (vm_title.west) -- node[above,font=\scriptsize]{1 \quad *} (user_title.east);
  \draw[arr] (cl_title.west) -- node[right,font=\scriptsize]{* \quad 1} ++(-2,0) |- (user_title.south);
\end{tikzpicture}
\caption{Modèle de données principal (entités ORM SQLAlchemy)}
\label{fig:class_diagram}
\end{figure}

% ─────────────────────────────────────────────────────────────────────────────
\section{Backlog Produit}
% ─────────────────────────────────────────────────────────────────────────────

Le tableau~\ref{tab:backlog} présente le backlog produit organisé par sprint.

\begin{table}[H]
\centering
\caption{Backlog produit}
\label{tab:backlog}
\begin{tabularx}{\linewidth}{clXl}
\toprule
\textbf{ID} & \textbf{Sprint} & \textbf{Histoire Utilisateur} & \textbf{Priorité} \\
\midrule
US-01 & 1 & S'inscrire et se connecter avec e-mail / mot de passe & Critique \\
US-02 & 1 & Configurer une VM via l'assistant multi-étapes & Critique \\
US-03 & 1 & Payer via le formulaire Stripe & Critique \\
US-04 & 1 & Provisionner la VM sur HCS via Terraform & Critique \\
US-05 & 1 & Diffuser la progression du provisionnement via SSE & Haute \\
US-06 & 1 & Recevoir un e-mail de confirmation après provisionnement & Moyenne \\
US-07 & 1 & Consulter et supprimer ses VM dans le tableau de bord & Critique \\
US-08 & 2 & Décrire sa charge de travail et recevoir une reco IA & Haute \\
US-09 & 2 & Parcourir les modèles de VM préconfigurés (marketplace) & Moyenne \\
US-10 & 2 & Ouvrir un terminal SSH navigateur vers sa VM & Haute \\
US-11 & 2 & Consulter le tableau de bord Netdata par VM & Haute \\
US-12 & 3 & Configurer et payer un cluster Kubernetes & Critique \\
US-13 & 3 & Provisionner le cluster K8s sur HCS avec kubeadm & Critique \\
US-14 & 3 & Diffuser la progression du provisionnement K8s via SSE & Haute \\
US-15 & 3 & Télécharger le kubeconfig depuis le tableau de bord & Critique \\
US-16 & 3 & Supprimer un cluster K8s depuis le tableau de bord & Haute \\
\bottomrule
\end{tabularx}
\end{table}

% ─────────────────────────────────────────────────────────────────────────────
\section{Architecture du Système}
% ─────────────────────────────────────────────────────────────────────────────

\subsection{Architecture Logique}

VM Marketplace adopte une \textbf{architecture client-serveur à trois services} avec
une séparation claire des responsabilités :

\begin{enumerate}
  \item \textbf{Frontend Next.js (port 3000)} --- rend l'interface utilisateur, gère
        l'état global via React Context, communique avec le backend via REST et SSE,
        et affiche un terminal WebSocket xterm.js.
  \item \textbf{Backend FastAPI (port 8000)} --- gère l'authentification, le CRUD des
        VM et clusters, le traitement des paiements, l'orchestration des sous-processus
        Terraform et l'installation Netdata via SSH.
  \item \textbf{Moteur IA (port 8001)} --- processus FastAPI distinct qui envoie des
        prompts en langage naturel à une instance Ollama locale, analyse la réponse
        JSON du LLM et la mappe vers une recommandation de flavor HCS.
\end{enumerate}

Le moteur IA est intentionnellement isolé du backend principal afin d'éviter de
bloquer l'API durant l'inférence LLM, qui peut prendre 10 à 30 secondes.

\subsection{Architecture de Déploiement}

La figure~\ref{fig:deploy_arch} présente l'architecture de déploiement avec toutes
les dépendances externes.

\begin{figure}[H]
\centering
\begin{tikzpicture}[
  box/.style={rectangle, draw, rounded corners=3pt, minimum width=3.5cm,
              minimum height=1cm, align=center, font=\small},
  cloud/.style={ellipse, draw, minimum width=2.8cm, minimum height=1cm,
                align=center, font=\small, fill=blue!10},
  arr/.style={-{Stealth}, thick},
  lbl/.style={font=\scriptsize, midway},
]
  % Serveur de développement
  \begin{scope}[local bounding box=dev]
    \node[box, fill=green!10]  (fe)     at (0,  3.5) {Frontend Next.js\\port 3000};
    \node[box, fill=orange!10] (be)     at (0,  2.0) {Backend FastAPI\\port 8000};
    \node[box, fill=purple!10] (ai)     at (0,  0.5) {Moteur IA\\port 8001};
    \node[box, fill=gray!10]   (ollama) at (0, -1.0) {Ollama\\port 11434\\llama3.1};
    \node[box, fill=yellow!10] (sqlite) at (0, -2.5) {SQLite\\vm\_marketplace.db};
    \node[box, fill=teal!10]   (tf)     at (0, -4.0) {Terraform\\huaweicloud/hcs};
  \end{scope}
  \node[draw, dashed, rounded corners, fit=(dev),
        label=above:{\small Serveur de développement}] {};

  % Services externes
  \node[cloud]                  (hcs)    at (6,  1.5) {Cloud HCS\\mesrscloud.rnu.tn};
  \node[cloud, fill=pink!20]    (stripe) at (6,  0.0) {API Stripe\\Paiement};
  \node[cloud, fill=cyan!10]    (resend) at (6, -1.5) {API Resend\\E-mail};

  % Ressources HCS
  \node[box, fill=blue!5] (vm)   at (10,  1.5) {Instances ECS\\+ EIP};
  \node[box, fill=blue!5] (k8s)  at (10,  0.0) {Cluster K8s\\Maître + Workers};
  \node[box, fill=blue!5] (netd) at (10, -1.5) {Netdata\\par VM};

  % Flèches
  \draw[arr] (fe)   -- node[lbl,above]{REST/SSE} (be);
  \draw[arr] (be)   -- node[lbl,right]{proxy}    (ai);
  \draw[arr] (ai)   -- (ollama);
  \draw[arr] (be)   -- (sqlite);
  \draw[arr] (be)   -- node[lbl,above]{sous-proc.} (tf);
  \draw[arr] (tf)   -- (hcs);
  \draw[arr] (be)   -- (stripe);
  \draw[arr] (be)   -- (resend);
  \draw[arr] (hcs)  -- (vm);
  \draw[arr] (hcs)  -- (k8s);
  \draw[arr] (be)   to[bend left=15] node[lbl,above]{SSH (Paramiko)} (vm);
  \draw[arr] (be)   to[bend left=20] node[lbl,above]{SSH (jump)} (k8s);
  \draw[arr] (vm)   -- (netd);
\end{tikzpicture}
\caption{Architecture de déploiement de VM Marketplace}
\label{fig:deploy_arch}
\end{figure}

% ─────────────────────────────────────────────────────────────────────────────
\section{Environnement de Travail}
% ─────────────────────────────────────────────────────────────────────────────

\begin{table}[H]
\centering
\caption{Environnement de développement}
\label{tab:env}
\begin{tabularx}{\linewidth}{lX}
\toprule
\textbf{Catégorie} & \textbf{Outils / Versions} \\
\midrule
Machine de développement & Windows 11 / macOS (projet multiplateforme) \\
IDE & Visual Studio Code avec extensions Python, ESLint, Prettier \\
Contrôle de version & Git + GitHub (dépôt \texttt{vm\_marketplace}) \\
Plateforme cloud & Huawei Cloud Stack --- \textit{mesrscloud.rnu.tn}, région \texttt{tn-global-1} \\
Terraform & v1.x avec fournisseur \texttt{huaweicloud/hcs} v2.4.0 \\
Python & 3.11 \\
Node.js & 18 LTS \\
Base de données & SQLite 3 (développement) ; fichier \texttt{vm\_marketplace.db} \\
Runtime LLM & Ollama v0.x avec le modèle \texttt{llama3.1} \\
\bottomrule
\end{tabularx}
\end{table}

% ─────────────────────────────────────────────────────────────────────────────
\section{Choix Technologiques}
% ─────────────────────────────────────────────────────────────────────────────

\subsection{Frontend : Next.js 14}

Next.js a été choisi pour son routage basé sur les fichiers, ses capacités de
rendu côté serveur et son écosystème mature. TypeScript garantit la sécurité des
types à toutes les frontières d'API. Tailwind CSS offre un style utilitaire rapide.
Stripe.js et la bibliothèque Stripe Elements gèrent la collecte des données de carte
conforme PCI dans le navigateur~\cite{nextjs_docs, tailwind_docs}.

\subsection{Backend : FastAPI}

FastAPI a été retenu pour sa gestion asynchrone des requêtes, sa documentation
OpenAPI automatique et ses annotations de types Python natives avec validation
Pydantic. Comparé à Flask (synchrone, moins structuré) et Django (couplage ORM plus
fort), FastAPI offre le meilleur équilibre entre performance et productivité pour un
service orienté API~\cite{fastapi_docs}. SQLAlchemy fournit la couche ORM pour SQLite.

\subsection{Infrastructure as Code : Terraform}

Terraform a été choisi pour le provisionnement des VM et des clusters car il fournit
des descriptions d'infrastructure déclaratives, une gestion d'état et des opérations
\texttt{apply} idempotentes. Le fournisseur \texttt{huaweicloud/hcs} (v2.4.0) supporte
tous les types de ressources HCS requis~\cite{hcs_terraform_provider}.

\subsection{Moteur IA : Ollama + Llama 3.1}

Le moteur de recommandation IA utilise Ollama pour servir le modèle Llama~3.1 (8B)
localement~\cite{ollama_docs, llama31}. Cela évite d'envoyer des descriptions de charge
de travail potentiellement sensibles à des API LLM cloud externes. Le moteur analyse
la sortie JSON structurée du LLM et la mappe vers un flavor HCS concret via un module
\texttt{mapper.py} personnalisé.

\subsection{Connectivité SSH : Paramiko}

Paramiko~\cite{paramiko_docs} est utilisé à deux fins : (1) l'installation
post-provisionnement de Netdata sur les VM individuelles via SSH direct ; (2) la
connexion SSH aux nœuds workers en IP privée via le maître comme hôte de rebond,
en utilisant \texttt{transport.open\_channel("direct-tcpip",...)}.

\subsection{Supervision : Netdata}

Netdata~\cite{netdata_docs} a été retenu pour son installation agent sans configuration
et son tableau de bord web intégré accessible sur le port 19999. Après installation,
l'URL Netdata de la VM est stockée en base de données et affichée sous forme d'iframe
dans la page de supervision de la plateforme.

% ─────────────────────────────────────────────────────────────────────────────
\section{Conclusion}
% ─────────────────────────────────────────────────────────────────────────────

Ce chapitre a formalisé les besoins, le modèle de données, le backlog produit et
l'architecture technique de VM Marketplace. L'architecture à trois services sépare
clairement le rendu de l'interface, la logique métier et l'inférence IA, tandis que
Terraform délègue toute la gestion des ressources HCS à un outil IaC déclaratif.
Le chapitre suivant présente les décisions de conception et d'implémentation prises
lors de chacun des trois sprints de développement.
